\documentclass[10pt,a4paper]{amsart}
\usepackage[margin=19mm]{geometry}
\usepackage{amsmath,amssymb,mathtools}
\usepackage[scr=rsfs,cal=euler]{mathalfa}
\usepackage{xfrac}
\usepackage[unicode,hidelinks,bookmarks=false]{hyperref}
\newcommand{\ie}{i.e.\ }
\newcommand{\cf}{cf.\ }
\newcommand{\resp}{respectively\ }
\newcommand{\Subsection}{\subsection}
\providecommand{\nobreakdash}{\nobreak\hbox{-}}
\setlength{\emergencystretch}{2em}
\multlinegap0pt
\usepackage[shortlabels]{enumitem}
\usepackage[normalem]{ulem}
\usepackage{xcolor}
\usepackage{amssymb}
\usepackage{mathtools}
\usepackage{algorithm}\usepackage{algpseudocode}
\usepackage{xfrac}
\usepackage{tikz}
\usepackage{tcolorbox}
\usepackage{float}
\newtheorem{assumption}{Assumption}
\newtheorem{lemma}{Lemma}
\newcommand{\E}{\mathbb E}
\newcommand{\R}{\mathbb R}
\newcommand{\beq}{\begin{equation}}
\newcommand{\bxi}{\boldsymbol{\xi}}
	\newcommand{\eeq}{\end{equation}}
\newcommand{\lln}{\left|}
\newcommand{\lnn}{\left\|}
\newcommand{\mr}[1]{{\color{blue}#1}}
\newcommand{\pc}[1]{{\color{cyan}#1}}
\newcommand{\rnn}{\right\|}  
\newcommand{\rrn}{\right|}
\newcommand{\us}[1]{\textcolor{red}{#1}}
\title{Complexity Guarantees for Zeroth-order Methods via Exponentially-shifted Gaussian Smoothing: Mitigating Dimension-dependence and Incorporating Decision-dependence}
\author{MINGRUI WANG, PRAKASH CHAKRABORTY, UDAY V. SHANBHAG}
\date{Source: arXiv:2604.15525v1 (CC BY 4.0)}
\begin{document}
\maketitle

\noindent\emph{Source and licence.} Complexity Guarantees for Zeroth-order Methods via Exponentially-shifted Gaussian Smoothing: Mitigating Dimension-dependence and Incorporating Decision-dependence. Version arXiv:2604.15525v1 (CC BY 4.0), distributed by arXiv under the Creative Commons Attribution 4.0 International licence, http://creativecommons.org/licenses/by/4.0/, as recorded in the arXiv licence metadata for this exact version and shown on the abstract page https://arxiv.org/abs/2604.15525v1. Full licence notice: \texttt{licenses/arxiv-2604.15525-CC-BY-4.0.txt}.\par\smallskip

\noindent\textit{Editorial scope.} Counted unit: the article's lemma lem:f-decdep1-Lip (source0.tex lines 1186--1189). The article's complete original proof of it is delivered with it (source0.tex lines 1190--1220, one \texttt{proof} environment of 31 lines). No mathematics is written, repaired or re-derived: the statement and the proof are the article's own lines, in the article's own notation, and no retained line is substituted or re-flowed.  Retained uncounted as the source-local context the proof needs: lines 546--551; lines 629--631; lines 648--650; lines 712--714; lines 758--773; lines 1009--1011; lines 1080--1092; lines 1119--1125; lines 1131--1133.  Nothing was omitted from any retained span, so the package carries no omission record.  No retained line contains a citation, so the wrapper carries no bibliography and no bibliography entry.  No retained line contains a figure environment, an image-inclusion command or drawing code, so no image is delivered and the package declares no asset dependency.  Editorial formatting only, in addition: an AMS \texttt{amsart} preamble that restates the article's own declarations and macros used by the retained lines, namely the environments \texttt{assumption}, \texttt{lemma} on separate counters, together with the standard packages those lines need.  The retained environments print in this document as Assumption 1, Lemma 1 in order of appearance, whereas the article prints the same items with the numbering of its own full document.  The complete native source of the article as distributed in the arXiv e-print for this exact version is bundled unchanged as \texttt{sources/198612/source0.tex}.
	\begin{assumption}\label{ass:f-Lip}
		$f: \R^n \mapsto \R$ is $L_0$-Lipschitz, i.e., there exists a constant $L_0>0$ such that
		$$
		\lln \, f(x) - f(y)\, \rrn \, \leq \,  L_0 \lnn \, x-y \, \rnn \text{ for all }x,y \in \R^n. 
		$$ 
	\end{assumption}

		\begin{equation}\label{eq:phieta}
			\phi_{\eta}(z) \, \triangleq \, \eta^{-n} \phi\left(\eta^{-1} z\right). 
		\end{equation}

		\begin{equation}\label{eq:feta}
			f_\eta(x) \, \triangleq \, f*\phi_\eta(x).
		\end{equation} 

			\begin{equation}\label{eq:GradientfetaD}
				\frac{\partial_i f_{\eta}(x)}{\partial x_i} \,  = \, \frac{1}{\eta \sqrt{2\pi}}\mathbb{E}_{V,Z^{-i}}\left[f\left(x_i+\eta\sqrt{2V},x^{-i}-Z^{-i}\right)-f\left(x_i-\eta\sqrt{2V},x^{-i}-Z^{-i}\right)\right].
			\end{equation}

		\begin{lemma}\label{Lemma:SmoothProperties} %\pc{Suggestion: we state part (a) elsewhere and state everything else for a generic function $f$ which is initially just $L_0$-Lipschitz. Then we use this lemma for both decision-dependent and not.}\\
			Assume $f$ satisfies Assumption~\ref{ass:f-Lip}, and $ \phi_{\eta}$, $f_{\eta}$ are defined as in \eqref{eq:phieta}, \eqref{eq:feta}, respectively. Then the following hold for any $\eta, \eta^{\prime} > 0$ and any $x,y \in \R^n$. %\mr{\bf MR: Shall we change the following space from $X$ to $\R^n$?}
			\begin{enumerate}[label=\alph*.]
				%\item \label{Lemma:a} %$f_{\eta} \in C^{\infty}$, and 
				%\sout{The components of $\nabla f_{\eta}$ are given by \eqref{eq:GradientfetaD}} \pc{(repetition of earlier Prop, lets get rid of this)}.
				\item \label{Lemma:b}$|f_{\eta}(x)-f_{\eta}(y)| \leq L_0\|x-y\|$ %for any $x, y \, \in \, \R^n$.
				\item \label{Lemma:c}$|f_{\eta}(x)-f(x)| \leq L_0\sqrt{n+1}\eta$ %for any $x \, \in \, \R^n$.
				\item \label{Lemma:d} If in addition $f$ is convex, then $f_\eta$ is convex and %for any $x \, \in \, \R^n$,
				\begin{equation}\label{eq:Lemmad}
					f(x) \leq f_\eta(x) \leq f(x)+L_0\sqrt{n+1}\eta.
				\end{equation}
				\item \label{Lemma:e} $\|f_\eta(x)-f_{\eta^{\prime}}(x)\| \leq L_0\sqrt{n+1}|\eta-\eta^{\prime}|$.
				\item \label{Lemma:f} $\|\nabla f_{\eta}(x) - \nabla f_{\eta}(y)\| \leq \frac{2L_0\sqrt{n}}{\eta \sqrt{2 \pi}}\|x-y\|$. %for any $x, y \, \in \, X$.
				\item \label{Lemma:g} If $f$ is differentiable and $L_1$-smooth, then $\|\nabla f_{\eta}(x) - \nabla f(x)\| \leq L_1\eta \sqrt{n}$. %for any $x \, \in \, X$.
			\end{enumerate}
		\end{lemma}

			\begin{equation}\label{eq:DDproblem}
				\min_{x \, \in \, X} f(x), \quad \mbox{ where } f(x) \, \triangleq  \, \mathbb{E}_{\bxi \sim \mathcal{D}(x)} \left[\, \hat{F}(x, \bxi)\, \right] \, = \, \int_{\Xi(x)} \hat{F}(x, \xi)p(\xi \mid  x) d\xi,
			\end{equation}

				\begin{assumption}\label{asm:decdep}
					For any $x,y\in\R^n$, the following hold 
					\begin{enumerate}
						\item $p(\xi \mid x)>0$ for any $\xi\in\Xi(x)$.
						\item\label{ass:p-x-y-dist} There exists an $L_\xi$ such that the {symmetric KL divergence} satisfies $D_{{\textnormal{KL}}}^{{\rm sym}}(p(\bullet \mid x) , p(\bullet \mid y)) \leq L_\xi^2\|x-y\|^2$. %\pc{(should we change $L_0$ to something else?)}
						\item\label{ass:p-ratio-bnd} For any $\xi\in\Xi(x)$, $\frac{p(\xi \mid x)}{\tilde{p}(\xi)}$ is uniformly bounded by $M$.
						\item\label{ass:f-hat-Lip} $\hat{F}(\bullet,\xi)$ is $\hat{L}_f$-Lipschitz continuous in the $L^2(\bxi \sim \tilde{p})$ metric, as specified next, i.e., for any $x, y \, \in \, \R^n$, the following holds.
						$$
						{\lnn \hat{F}(x, \bullet) - \hat{F}(y, \bullet) \rnn}_{L^2(\bxi \sim \tilde{p})} \, \leq \, \hat{L}_f \|x-y\|.
						$$
						Furthermore, $\hat{F}(\bullet,\xi)$ is bounded by $M_f$ for almost every $\xi\in\Xi(\bullet)$. 
					\end{enumerate}
				\end{assumption}

				\begin{align}\label{eq:seclast}
					\notag     \E_{\bxi \sim \tilde{p}} &\left[{\lln \dfrac{p(\bxi | x)}{\tilde{p}(\bxi)} - \dfrac{p(\bxi | y)}{\tilde{p}(\bxi)} \rrn}^2 \right] 
					= 
					\int \left|\frac{p(\xi \mid x)}{\tilde{p}(\xi)}-\frac{p(\xi \mid y)}{\tilde{p}(\xi)}\right|^2 \tilde{p}(\xi)d\xi=\int \frac{\left|p(\xi \mid x)-p(\xi \mid y)\right|^2}{\tilde{p}(\xi)} d\xi\nonumber\\
					&=\int \frac{\left(p(\xi \mid x)-p(\xi \mid y)\right)_+^2}{p(\xi \mid x)}\frac{p(\xi \mid x)}{\tilde{p}(\xi)} d\xi+\int \frac{\left(p(\xi \mid x)-p(\xi \mid y)\right)_-^2}{p(\xi \mid y)}\frac{p(\xi \mid y)}{\tilde{p}(\xi)} d\xi\nonumber\\
					& \leq Md^2_2\left(p(\xi \mid x),p(\xi \mid y)\right)+Md^2_2\left(p(\xi \mid y),p(\xi \mid x)\right),
				\end{align}

				\begin{equation}\label{eq:KBbound}
					d^2_2\left(p(\xi \mid x),p(\xi \mid y)\right)+d^2_2\left(p(\xi \mid y),p(\xi \mid x)\right) \leq D_{{\textnormal{KL}}}^{{\rm sym}}\left(p(\xi \mid x),p(\xi \mid y)\right) \leq L_\xi^2{\|x-y\|}^2
				\end{equation}

						\begin{lemma}\label{lem:f-decdep1-Lip}
							Under Assumption~\ref{asm:decdep}, $f$ defined in \eqref{eq:DDproblem} is $\hat{L}_0$-Lipschitz continuous, where $\hat{L}_0=M \hat{L}_f+M_fL_\xi$. Consequently, Lemma~\ref{Lemma:SmoothProperties} holds true with Lipschitz constant $\hat{L}_0$. %all previous results we derived for non-decision-dependent problem will hold for decision-dependent problem. 
							%\us{This needs clarification. We may want to relate to Lemma 3.2.}
						\end{lemma}

						\begin{proof}
						Recall from \eqref{eq:DDproblem} that $f(x)=\int_{\Xi} \hat{F}(x, \xi)p(\xi \mid  x) d\xi$. Consequently
						\begin{align}\label{eq:fdecLip1}
							|\, f(x)-f(y)\, |&=\left|\int_{\Xi} \hat{F}(x, \xi)p(\xi \mid  x) d\xi-\int_{\Xi} \hat{F}(y, \xi)p(\xi \mid  y) d\xi\right|\nonumber\\
							& \leq \int_\Xi\left|\hat{F}(x,\xi)-\hat{F}(y,\xi)\right| p(\xi\mid x) d\xi+\int_\Xi\left|\hat{F}(y,\xi)\right|\left|p(\xi \mid x)-p(\xi \mid y)\right|d\xi.
							%&\leq \mr{L_\xi}{\|x-y\|}+\mr{M_f}\int_\Xi \left|p(\xi \mid x)-p(\xi \mid y)\right|d\xi
						\end{align}
						By Lipschitz continuity of $\hat{F}(\bullet,\xi)$ and the boundedness of $\frac{p(\xi\mid x)}{\tilde{p}(\xi)}$ from Assumption~\ref{asm:decdep}, we may bound the first term in \eqref{eq:fdecLip1} as follows.
						\beq\label{eq:fdecLip-a}
						\int_\Xi\left|\hat{F}(x,\xi)-\hat{F}(y,\xi)\right| p(\xi\mid x) d\xi=\int_\Xi\left|\hat{F}(x,\xi)-\hat{F}(y,\xi)\right|\tilde{p}(\xi) \frac{p(\xi\mid x)}{\tilde{p}(\xi)} d\xi \, \leq \,  \hat{L}_f M  {\|x-y\|}.
						\eeq
						By Pinsker's Theorem, the total variation distance $D_{\textnormal{TV}}(p(\bullet|x), p(\bullet|y)) =\frac{1}{2} \int |p(\xi|x) - p(\xi|y)| d\xi \leq \sqrt{\frac{1}{2} D_{\textnormal{KL}}(p(\bullet|x), p(\bullet|y))}$. Consequently by the elementary inequality $\sqrt{a/2}+\sqrt{b/2}<\sqrt{a+b}$ and Assumption~\ref{asm:decdep}.\ref{ass:p-x-y-dist}, we obtain
						\beq\label{eq:fdecLip-c}
						\int \lln p(\xi | x) - p(\xi|y) \rrn d\xi \leq \sqrt{D_{{\textnormal{KL}}}^{{\rm sym}}(p(\bullet |x), p(\bullet | y))} \leq L_\xi \|x-y\|. 
						\eeq
						%\us{\bf UVS: I'm sorry I missed where we used $\sqrt{a/2}+\sqrt{b/2}<\sqrt{a+b}$} \pc{(D = symmetric-KL)}
						By boundedness of $\hat{F}{(\bullet,\xi)}$ from Assumption~\ref{asm:decdep}.\ref{ass:f-hat-Lip} and \eqref{eq:fdecLip-c}, we obtain that
						\beq\label{eq:fdecLip-b}
						\int_\Xi\left|\hat{F}(y,\xi)\right|\left|p(\xi \mid x)-p(\xi \mid y)\right|d\xi \leq M_f \int \lln p(\xi | x) - p(\xi|y) \rrn d\xi \leq M_f L_\xi {\|x-y\|}.
						\eeq
						Plugging the inequalities \eqref{eq:fdecLip-a} and \eqref{eq:fdecLip-b} in the relation \eqref{eq:fdecLip1} we obtain our desired result.
						%\us{By assumption,} we know that there exists \us{a density}  $\tilde{p}$,
						%\begin{equation*}
						%	\int_\Xi \left|p(\xi \mid x)-p(\xi \mid y)\right|d\xi = \int_\Xi \left|\frac{p(\xi \mid x)-p(\xi \mid y)}{\tilde{p}(\xi)}\right|\tilde{p}(\xi)d\xi \leq \left(\int_\Xi \left|\frac{p(\xi \mid x)-p(\xi \mid y)}{\tilde{p}(\xi)}\right|^2\tilde{p}(\xi)d\xi\right)^{1/2},
						%\end{equation*}
						%\pc{where the last inequality follows from Cauchy-Schwarz inequality. Then following the same arguments as in} \eqref{eq:seclast} and \eqref{eq:KBbound} we get
						%\begin{equation*}
						%	\int_\Xi \left|p(\xi \mid x)-p(\xi \mid y)\right|d\xi \leq \mr{L_\xi}{\|x-y\|}.
						%\end{equation*}
						%\pc{Plugging} in the above inequality to \eqref{eq:fdecLip1} we get the Lipschitz property.
						\end{proof}

\end{document}
